\documentclass[10pt, aspectratio=169]{beamer}
\input{../../_en_preambulo_beamer}
% Comandos y macros estadísticas compartidas para las presentaciones Beamer en inglés (Metropolis)

% Conjuntos numéricos básicos
\newcommand{\R}{\mathbb{R}}
\newcommand{\N}{\mathbb{N}}
\newcommand{\Q}{\mathbb{Q}}
\newcommand{\Z}{\mathbb{Z}}
\newcommand{\set}[1]{\left\{ #1 \right\}}
\newcommand{\abs}[1]{\left|#1\right|}
\newcommand{\norm}[1]{\left\| #1 \right\|}

% Comandos estadísticos y probabilísticos del curso
\newcommand{\Var}{\operatorname{Var}}
\newcommand{\cov}{\operatorname{Cov}}
\newcommand{\corr}{\rho}
\newcommand{\comb}[2]{\begin{pmatrix} #1 \\ #2 \end{pmatrix}}
\newcommand{\s}{\sigma}
\newcommand{\particion}{\mathfrak{P}}
\newcommand{\card}[1]{\#\left( #1 \right)}
\newcommand{\seq}[3]{\set{#1}_{#2}^{#3}}
\newcommand{\E}{\mathbb{E}}
\newcommand{\Prob}{\mathbb{P}}

% Entornos pedagógicos y cajas en inglés académico natural (soportando tanto nombres en inglés como en español)
\newenvironment{discussion}{%
  \begin{alertblock}{Discussion Question}%
}{%
  \end{alertblock}%
}
\newenvironment{discusion}{%
  \begin{alertblock}{Discussion Question}%
}{%
  \end{alertblock}%
}

\newenvironment{reflection}{%
  \begin{alertblock}{Classroom Reflection}%
}{%
  \end{alertblock}%
}
\newenvironment{reflexion}{%
  \begin{alertblock}{Classroom Reflection}%
}{%
  \end{alertblock}%
}

\newenvironment{paradox}{%
  \begin{alertblock}{Exploratory Paradox}%
}{%
  \end{alertblock}%
}
\newenvironment{paradoja}{%
  \begin{alertblock}{Exploratory Paradox}%
}{%
  \end{alertblock}%
}

\newenvironment{motivation}{%
  \begin{exampleblock}{Motivation: Data Science in Practice}%
}{%
  \end{exampleblock}%
}
\newenvironment{motivacion}{%
  \begin{exampleblock}{Motivation: Data Science in Practice}%
}{%
  \end{exampleblock}%
}

\newenvironment{quickchallenge}{%
  \begin{block}{Quick Team Challenge}%
}{%
  \end{block}%
}
\newenvironment{retoclase}{%
  \begin{block}{Quick Team Challenge}%
}{%
  \end{block}%
}


\title[Regression Diagnostics]{Section 09.09: Residual Diagnostics and Multicollinearity}
\subtitle{Classical Assumptions, VIF, and Influential Observations}
\author[J. Castillo Colmenares]{Juliho Castillo Colmenares}
\institute{Tecnológico de Monterrey}
\date{\vspace{-1.5cm}}

\begin{document}

% Slide 1: Title
\begin{frame}[plain]
  \titlepage
\end{frame}

% Slide 2: Roadmap
\begin{frame}{Roadmap --- Chapter 09: Linear and Multiple Regressions}
  \scriptsize
  Where are we in the modular development of this chapter?
  \begin{itemize}\itemsep=0.02em
    \item \textbf{09.01} Correlation as the Premise of Regression
    \item \textbf{09.02} Introduction to Linear Regression
    \item \textbf{09.03} The Mathematics of Regression: OLS Derivation and $R^2$
    \item \textbf{09.04} Linear Regression on Simulated Data
    \item \textbf{09.05} Optimal Coefficients, $t$/$F$ Tests, and RSE
    \item \textbf{09.06} Implementation in Python with \texttt{statsmodels}
    \item \textbf{09.07} Multiple Regression, the Normal Equation, and Ridge/Lasso
    \item \textbf{09.08} Model Validation and $k$-fold Cross-Validation
    \item \textbf{09.09} Residual Diagnostics and Multicollinearity (VIF) \emph{(Today)}
    \item \textbf{09.10} Regression with \texttt{scikit-learn}
    \item \textbf{09.11} Categorical Variables and Dummy Variables
    \item \textbf{09.12} Nonlinear Transformations and Polynomial Regression
  \end{itemize}
  \pause
  \vspace{-0.05cm}
  \begin{block}{Session Objective}
    Rigorously audit whether a fitted regression model is \emph{trustworthy}: verify the classical assumptions via residual analysis and formal tests (Durbin-Watson, Breusch-Pagan), quantify multicollinearity with the Variance Inflation Factor, and identify influential observations via Cook's Distance.
  \end{block}
\end{frame}

% Slide 3: Motivation
\begin{frame}{Is the Model We Fit Trustworthy?}
  \footnotesize
  \begin{motivacion}
    In Sections 09.03 through 09.07 we built and solved the regression model (simple and multiple), but all inference (intervals, $t$/$F$ tests) rests on assumptions about the error term $\varepsilon$ that are \textbf{never automatically verified}. A model with high $R^2$ can be completely invalid if its residuals violate these assumptions.
  \end{motivacion}

  \vspace{0.15cm}
  \pause
  \begin{itemize}\itemsep=0.1cm
    \item \textbf{Residuals} $e_i=Y_i-\hat Y_i$ are the primary diagnostic tool. \pause
    \item We complement graphical analysis with formal tests and the identification of problematic observations.
  \end{itemize}
\end{frame}

% Slide 4: Classical assumptions and residuals
\begin{frame}{Classical Assumptions and Residual Analysis}
  \footnotesize
  \begin{block}{Five Assumptions on $\varepsilon_i$}
    Linearity, Independence, Homoscedasticity ($\text{Var}(\varepsilon_i)=\sigma^2$), Normality ($\varepsilon_i\sim N(0,\sigma^2)$), and Zero Mean.
  \end{block}
  \pause

  \vspace{0.1cm}
  \begin{alertblock}{Patterns in the Residuals-vs-Fitted Plot}
    U-shape $\to$ non-linearity; \textbf{funnel shape} $\to$ heteroscedasticity; isolated extreme points $\to$ possible outliers.
  \end{alertblock}
\end{frame}

% Slide 5: Formal tests
\begin{frame}{Formal Tests: Durbin-Watson and Breusch-Pagan}
  \footnotesize
  \begin{block}{Durbin-Watson (Independence)}
    $$DW=\frac{\sum_{i=2}^n(e_i-e_{i-1})^2}{\sum_{i=1}^n e_i^2}\approx2(1-\hat\rho), \qquad DW\in(0,4),\ DW\approx2 \text{ ideal}$$
  \end{block}
  \pause

  \vspace{0.1cm}
  \begin{alertblock}{Breusch-Pagan (Homoscedasticity)}
    Auxiliary regression of $\hat\varepsilon_i^2$ on the predictors: $$LM=n\cdot R_{\text{aux}}^2 \sim \chi^2_p \text{ under } H_0\text{: homoscedasticity}$$
  \end{alertblock}
\end{frame}

% Slide 6: VIF
\begin{frame}{Multicollinearity: the Variance Inflation Factor}
  \footnotesize
  \begin{block}{Operational Definition}
    $$\text{VIF}_j=\frac{1}{1-R_j^2}, \qquad R_j^2: \text{ coeff. of regressing } X_j \text{ on the other predictors}$$
  \end{block}
  \pause

  \vspace{0.1cm}
  \begin{alertblock}{Identity with the Correlation Matrix}
    $\text{VIF}_j=(\mathbf{R}^{-1})_{jj}$: a high VIF is the diagnostic manifestation that $\mathbf{X}^T\mathbf{X}$ is approaching singularity (Section 09.07). Conventional threshold: $\text{VIF}>5$ warrants Ridge, Lasso, or variable elimination.
  \end{alertblock}
\end{frame}

% Slide 7: Influential observations
\begin{frame}{Influential Observations: Cook's Distance}
  \footnotesize
  \begin{block}{Computational Formula}
    $$D_i=\frac{r_i^2}{p+1}\cdot\frac{h_{ii}}{1-h_{ii}}, \qquad \text{threshold: } D_i>4/n$$
  \end{block}
  \pause

  \vspace{0.1cm}
  \begin{alertblock}{Two Distinct Mechanisms}
    High leverage $h_{ii}$ (extreme position in $X$) \textbf{or} a large residual $r_i$ (possible outlier in $Y$) --- either one, or both, can drive $D_i$ up.
  \end{alertblock}
\end{frame}

% Slide 8: Python Lab (Block 1)
\begin{frame}[fragile]{Python Lab: Durbin-Watson and Breusch-Pagan}
  \scriptsize
  Both tests implemented from scratch with \texttt{numpy}/\texttt{scipy} (\texttt{09.03\_regression\_diagnostics.py}):
  {\lstset{basicstyle=\fontsize{3.3pt}{4.0pt}\selectfont\ttfamily,aboveskip=0.05em,belowskip=0.05em}
  \lstinputlisting[firstline=19, lastline=41]{../../code/09_regresiones/09.09_regression_diagnostics.py}}
\end{frame}

% Slide 9: Python Lab (Block 2)
\begin{frame}[fragile]{Python Lab: VIF and its Matrix Identity}
  \scriptsize
  VIF computed via auxiliary regression and verified against $\text{diag}(\mathbf{R}^{-1})$:
  {\lstset{basicstyle=\fontsize{3.4pt}{4.1pt}\selectfont\ttfamily,aboveskip=0.05em,belowskip=0.05em}
  \lstinputlisting[firstline=44, lastline=70]{../../code/09_regresiones/09.09_regression_diagnostics.py}}
\end{frame}

% Slide 10: Python Lab (Block 3)
\begin{frame}[fragile]{Python Lab: Cook's Distance}
  \scriptsize
  Detecting influential observations via high leverage vs. large residual:
  {\lstset{basicstyle=\fontsize{3.2pt}{3.9pt}\selectfont\ttfamily,aboveskip=0.05em,belowskip=0.05em}
  \lstinputlisting[firstline=73, lastline=103]{../../code/09_regresiones/09.09_regression_diagnostics.py}}
\end{frame}

% Slide 11: Terminal Output
\begin{frame}[fragile]{Python Lab: Terminal Output}
  \scriptsize
  Standard output produced when running the lab script:
  \vspace{0.02cm}
  \begin{block}{Standard Console Output}
    \fontsize{3.6pt}{4.3pt}\selectfont
    \begin{verbatim}
=== Block 1: Durbin-Watson and Breusch-Pagan ===
DW=0.5759 (approx 2(1-rho)=0.7866) -> strong positive autocorrelation
Breusch-Pagan: R2_aux=0.1760, LM=10.5576, p=0.0051 -> heteroscedasticity

=== Block 2: VIF ===
Predictor 1: VIF=94.59  Predictor 2: VIF=94.55  Predictor 3: VIF=1.01
VIF via 1/(1-R^2) matches VIF via diag(R^-1) exactly

=== Block 3: Cook's Distance ===
Obs 0: leverage=0.59, resid=-0.85, D=0.52 -> INFLUENTIAL (leverage-driven)
Obs 1: leverage=0.04, resid=5.10,  D=0.54 -> INFLUENTIAL (residual-driven)
    \end{verbatim}
  \end{block}
\end{frame}


% Slide 16: Closing
\begin{frame}{Section Closing: Regression Diagnostics}
  \begin{reflexion}
    A fitted model is never trustworthy by decree: we have audited the classical assumptions via residuals and formal tests (Durbin-Watson, Breusch-Pagan), quantified multicollinearity with VIF, and detected influential observations via Cook's Distance --- closing the full cycle of internal model validation.
  \end{reflexion}

  \vspace{0.2cm}
  \pause
  \begin{block}{Key Takeaways}
    \begin{itemize}\itemsep=0.08cm
      \item \textbf{$DW\approx2$:} the ideal value; deviations indicate positive or negative autocorrelation.
      \item \textbf{VIF$_j=(\mathbf{R}^{-1})_{jj}$:} a high VIF reflects the same numerical instability of $\mathbf{X}^T\mathbf{X}$ seen in Section 09.07.
      \item \textbf{Cook's Distance:} high leverage and large residuals are independent mechanisms of influence.
    \end{itemize}
  \end{block}
\end{frame}

% Slide 17: Modular Perspective
\begin{frame}{Modular Perspective --- The Next Challenge}
  \scriptsize
  \begin{retoclase}
    With the model internally audited and already externally validated in Section 09.08, we close the chapter's statistical core. The three final sections (09.10-09.12) are more applied in nature: Section 09.10 shows how everything built by hand in \texttt{numpy}/\texttt{scipy} translates into professional practice with \texttt{scikit-learn}.
  \end{retoclase}

  \vspace{0.08cm}
  \pause
  \begin{alertblock}{Recommended Preparation}
    \begin{itemize}\itemsep=0.06cm
      \item Reflect on why a model can pass every diagnostic in this section and still generalize poorly (overfitting) --- that's why the validation in Section 09.08 is independent of this internal diagnosis.
      \item Review the basic \texttt{scikit-learn} API (\texttt{fit}, \texttt{predict}, \texttt{score}) that unifies the interface across all of the library's models.
    \end{itemize}
  \end{alertblock}
\end{frame}

\end{document}
